\section{Flip classes of ordinary graphs}\label{sec:ordinary}

In this section we prove Theorem~\ref{thm:main} for ordinary graphs, assuming a
connectivity statement, Theorem~\ref{thm:connectivity}, which we prove in
Section~\ref{sec:connectivity}.

The idea is as follows.
The four-term identity gives linear relations among the sums $\Kv(\Lay_s)$ and
$\Ev(\Lay_s)$ over all orientations with $s$ strict and ascending edges.
These are the bounce relations of~\cite{AlexanderssonSulzgruber2022} in a form
that keeps the number of edges fixed.
To pass from all orientations to one flip class $C$, we show that every other
class with the same number of edges can be written, with the same integer coefficients for
$\Kv$ and $\Ev$, in terms of classes of smaller graphs.
For this we need to know that one class contains all orientations of a certain
kind; this is Theorem~\ref{thm:connectivity}.
The relation is first obtained for an initial piece of the graph, called the
prefix, and then transported to the whole graph by fixing the color of the last
vertex of the prefix.
Throughout, $\Gamma$ is an ordinary decorated unit-interval graph on $[n]$, and
$|E(\Gamma)|=\sum_ia_i$ denotes its number of natural edges.

\subsection{Changing the decoration}

The values $\Kv$ and $\Ev$ only depend on the posets $\Pplus{\theta}$, and the
flip relation only depends on the graphs $\theta\cup S$.
Hence the following two operations transport flip classes to flip classes and
preserve both values.

\begin{lemma}\label{lem:decoration}
Let $C$ be a flip class of $\Gamma$ and let $f$ be a corner.
\begin{enumerate}[(i)]
\item If $f$ is an edge of $\theta\cup S$ for all $\theta\in C$, then marking
$f$ strict, or unmarking it if it is strict, gives a graph $\Gamma'$ in which
$C$ is again a flip class, with the same graphs $\theta\cup S$.
\item If $f$ is an edge of $\theta\cup S$ for no $\theta\in C$, then deleting
$f$ gives a decorated unit-interval graph $\Gamma\setminus f$ with
$|E(\Gamma)|-1$ natural edges, in which $C$ is again a flip class.
\end{enumerate}
\end{lemma}

\begin{proof}
A corner is never a short edge of a flip.
In (i) we therefore only change the label of an edge which is present
throughout, and the flips are the same.
In (ii), deleting the maximal interval $f$ keeps the natural edges closed under
subintervals, and the remaining strict edges are still corners.
No orientation in $C$ uses $f$, so $f$ is not the long edge of a flip from $C$,
and conversely every flip of $\Gamma\setminus f$ is a flip of $\Gamma$.
\end{proof}

By Lemma~\ref{lem:flipInvariants}, every corner is either present throughout
a flip class or absent throughout.

\subsection{Fixing the color of a vertex}

We shall need the two values with the color of one vertex $c$ fixed.
For an orientation $\theta$ and $\beta\in\setNN$, let
$\Kp_\theta[\beta]$ be the sum of $x_\kappa$ over the strict
$\Pplus{\theta}$-partitions $\kappa$ with $\kappa(c)=\beta$.
For the elementary value, replace every fibre of $\hrv{\theta}{\cdot}$ by the
chain on its elements in increasing order.
The strict partitions of this disjoint union of chains give
$\elementaryE_{\lambda(\theta)}$, and we let $\elementaryE_\theta[\beta]$ be the sum over
those with $\kappa(c)=\beta$.
If $c$ is the largest vertex of its fibre, which has $p$ elements, then
\begin{equation}\label{eq:markedE}
 \elementaryE_\theta[\beta]
 =x_\beta\,\elementaryE_{p-1}(x_1,\dotsc,x_{\beta-1})
  \prod_{B}\elementaryE_{|B|}(\xvec),
\end{equation}
the product ranging over the other fibres $B$.
For a set $\Omega$ of orientations we write $\Kv(\Omega)[\beta]$ and $\Ev(\Omega)[\beta]$ for the
corresponding sums.
Summing over $\beta$ gives back $\Kv(\Omega)$ and $\Ev(\Omega)$.

Below, $c$ is always the largest vertex of the graph under consideration, so
that~\eqref{eq:markedE} applies, and $\elementaryE_\theta[\beta]$ only depends on
the fibre sizes and on the size of the fibre of $c$.
The values with the color of $c$ fixed are preserved by every operation used
below: by
Lemma~\ref{lem:decoration}; by deleting a redundant edge, which preserves
$\Pplus{\theta}$; by the four-term identity, whose pairing preserves the whole
map $\hrv{\theta}{\cdot}$; and by relabellings as in Lemma~\ref{lem:sinkRelabel}
which fix $c$.

\subsection{Two relations}

For an integer $s$, let $\defin{\Lay_s(\Gamma)}$ be the set of orientations
$\theta$ such that $\theta\cup S$ has $s$ edges.
By Lemma~\ref{lem:flipInvariants}, $\Lay_s(\Gamma)$ is a union of flip
classes.
We use two local configurations, see Figure~\ref{fig:contacts}.

\begin{definition}\label{def:contacts}
Let $x<y<t$ be vertices with $y=x+1$, such that $xt$ is a strict edge, $y$ is not
the endpoint of any strict edge, and $xt$ is the only strict edge leaving $x$.
Assume that no natural edge joins a vertex at most $y$ to a vertex larger than
$t$, that is, $b(t+1)>y$, where we put $b(n+1)=n+1$.
\begin{enumerate}[(i)]
\item The triple $(x,y,t)$ is an \defin{NN contact} if $b(x)=b(y)$ and no strict
edge ends at $x$.
Its \defin{target} $R$ is obtained by deleting the edge $xt$ and marking the
corner $yt$ strict.
\item Let $u<x$. The quadruple $(u,x,y,t)$ is a \defin{DN contact} if $ux$ is a
strict edge, no other strict edge ends at $x$, and $b(y)=u+1$.
Its \defin{target} $R$ is obtained by changing $a_y$ to $a_y+1$ and $a_t$ to
$a_t-1$, so that $uy$ becomes a natural edge and $xt$ does not, and by replacing
the strict edges $ux,xt$ by $uy,yt$.
\end{enumerate}
\end{definition}

One checks directly from the inequalities that $R$ is a decorated
unit-interval graph.
In the DN case, $a_y=x-u=a_x$, so the new area sequence is valid at $y$, and
$a_{t+1}<a_t$ by the condition on $b(t+1)$.
The target of an NN contact has one natural edge fewer.
The target of a DN contact has the same number of natural edges, and
$\mom(\Gamma)\coloneqq\sum_i i\,a_i$ is smaller by $t-y$.

\begin{figure}[!ht]
\centering
\begin{tikzpicture}[llt={cell=0.8em}]
\LLTframe{5}{8}
\draw[llt arbitrary] (0,0)--(1,2) (1,4)--(2,6) (4,7)--(5,8);
\LLTsteps{1}{2}{nn}\LLTsteps{2}{6}{de}
\draw[llt bounce] (3,7)--(3,3)--(1,3);
\end{tikzpicture}
\;$\longmapsto$\;
\begin{tikzpicture}[llt={cell=0.8em}]
\LLTframe{5}{8}
\draw[llt arbitrary] (0,0)--(1,2) (1,4)--(2,6) (4,7)--(5,8);
\LLTsteps{1}{2}{nn}\LLTsteps{2}{6}{ed}
\draw[llt bounce] (3,6)--(3,3)--(1,3);
\end{tikzpicture}
\qquad\qquad
\begin{tikzpicture}[llt={cell=0.8em}]
\LLTframe{5}{8}
\draw[llt arbitrary] (0,0)--(1,2) (2,4)--(2,6) (4,7)--(5,8);
\LLTsteps{1}{2}{dn}\LLTsteps{2}{6}{de}
\draw[llt bounce] (3,7)--(3,3)--(2,3);
\end{tikzpicture}
\;$\longmapsto$\;
\begin{tikzpicture}[llt={cell=0.8em}]
\LLTframe{5}{8}
\draw[llt arbitrary] (0,0)--(1,2) (2,4)--(2,6) (4,7)--(5,8);
\LLTsteps{1}{2}{nd}\LLTsteps{2}{6}{ed}
\draw[llt bounce] (3,6)--(3,3)--(1,3);
\end{tikzpicture}
\caption{An NN contact (left) and a DN contact (right), each followed by its
target, drawn as in~\cite[(3), (4)]{AlexanderssonSulzgruber2022}.
Thick segments are steps of the Schröder path, and dotted segments are
arbitrary.
The diagonal steps at the top cross the strict edge $xt$, and $yt$ in the
targets.
The dashed red path goes down from there to the diagonal box of $x$ and then left to
the path, between the rows of $x$ and $y=x+1$.
In the DN contact the diagonal step at the bottom crosses the strict edge $ux$,
and in its target $uy$.}
\label{fig:contacts}
\end{figure}

\begin{proposition}\label{prop:layer}
Let $\Psi$ be one of $\Kv$, $\Ev$, $\Kv[\beta]$ or $\Ev[\beta]$, with the color
fixed at $c=n$.
If $\Gamma$ has an NN contact with target $R$, then for all $s$
\[
 \Psi(\Lay_s(\Gamma))=\Psi(\Lay_s(R))+\Psi(\Lay_{s-1}(R)).
\]
If $\Gamma$ has a DN contact with target $R$, then $\Psi(\Lay_s(\Gamma))=\Psi(\Lay_s(R))$
for all $s$.
\end{proposition}

\begin{proof}
For an orientation $\theta$ we write $\Psi(\theta)=\Psi(\{\theta\})$.
Comparing coefficients of $q^s$, the two relations are equivalent to
\begin{equation}\label{eq:layerGenerating}
 \sum_{\theta\in\ORIENTS(\Gamma)}q^{|\theta\cup S|}\,\Psi(\theta)
 =(1+q)\sum_{\theta\in\ORIENTS(R)}q^{|\theta\cup S|}\,\Psi(\theta)
\end{equation}
for an NN contact, and to the same identity without the factor $1+q$ for a DN
contact; here $q$ is a formal variable.

Consider an NN contact $(x,y,t)$, and let $\sigma$ be the transposition of $x$
and $y$.
Let $E'$ be the set of natural edges other than $xy$, $yt$ and $xt$.
It is the same for $\Gamma$ and $R$, and $\sigma$ permutes it: for $p<x$, the
pair $px$ is natural if and only if $py$ is, since $b(x)=b(y)$; for $y<p<t$, the
pairs $xp$ and $yp$ are both natural since they lie inside $xt$; and there are no
natural edges from $x$ or $y$ to vertices beyond $t$.
The strict edges in $E'$ avoid $x$ and $y$, so $\sigma$ fixes them.

We group the orientations $\theta$ of $\Gamma$ and of $R$ by the set
$\eta=(\theta\cup S)\cap E'$.
In the pictures below, the vertices $x<y<t$ are drawn in increasing order, the
dotted line stands for $\eta$, thick red arrows are strict edges and thin black
arrows are ascending edges, and each picture stands for the value $\Psi$ of the
corresponding set of edges.
For fixed $\eta$, the orientations of $\Gamma$ contribute $q^{|\eta|}$ times
\[
 q\,\EdgePic{3}{x,y,t}{1/3/s}
 +q^2\,\EdgePic{3}{x,y,t}{1/3/s,1/2/a}
 +q^2\,\EdgePic{3}{x,y,t}{1/3/s,2/3/a}
 +q^3\,\EdgePic{3}{x,y,t}{1/3/s,1/2/a,2/3/a}\,,
\]
and the orientations of $R$ contribute $q^{|\eta|}$ times
\[
 q\,\EdgePic{3}{x,y,t}{2/3/s}
 +q^2\,\EdgePic{3}{x,y,t}{2/3/s,1/2/a}\,.
\]
By the four-term identity, and since $xt$ is redundant once $xy$ and $yt$ are
present,
\[
 \EdgePic{3}{x,y,t}{1/3/s,1/2/a}
 +\EdgePic{3}{x,y,t}{1/3/s,2/3/a}
 =\EdgePic{3}{x,y,t}{1/3/s}
 +\EdgePic{3}{x,y,t}{2/3/s,1/2/a}\,.
\]
Hence the orientations of $\Gamma$ with fixed $\eta$ contribute $(1+q)\,q^{|\eta|}$
times
\[
 q\,\EdgePic{3}{x,y,t}{1/3/s}
 +q^2\,\EdgePic{3}{x,y,t}{2/3/s,1/2/a}\,.
\]
Finally, $\sigma$ maps $\eta+xt$ onto $\sigma(\eta)+yt$.
Its only inversion is the pair $x<y$, and $x$ is not a sink of $\eta+xt$, so
Lemma~\ref{lem:sinkRelabel} applies; $\sigma$ fixes $c=n$, and relabelling
colorings by $\sigma$ preserves $\Kp$.
Hence the first picture for $\eta$ has the same value as the first picture of
$R$ for $\sigma(\eta)$.
Since $\sigma$ permutes the possible sets $\eta$ and preserves their sizes,
summing over $\eta$ gives~\eqref{eq:layerGenerating}.

For a DN contact $(u,x,y,t)$, let $E'$ be the set of natural edges of $\Gamma$ or
$R$ other than $ux$, $uy$, $xy$, $xt$ and $yt$, and again let $\sigma$ exchange
$x$ and $y$.
The previous argument shows that $\sigma$ permutes $E'$, now using that for
$p<x$ with $p\neq u$, the pair $px$ is natural if and only if $py$ is.
With the vertices $u<x<y<t$, the source has the strict edges $ux,xt$ and the
target the strict edges $uy,yt$, and the two remaining local edges of each may be
ascending or not.
We apply the four-term identity to the triangle $x<y<t$ in the source and to the
triangle $u<x<y$ in the target.
It remains to compare the terms with no ascending local edge, and with two:
\begin{gather*}
 \EdgePic{4}{u,x,y,t}{1/2/s,2/4/s}
 \;\longleftrightarrow\;
 \EdgePic{4}{u,x,y,t}{1/3/s,3/4/s}
 \\[0.6ex]
 \EdgePic{4}{u,x,y,t}{1/2/s,2/4/s,2/3/a,3/4/a}
 =\EdgePic{4}{u,x,y,t}{1/3/s,3/4/s,1/2/a,2/3/a}\,.
\end{gather*}
In the first line, $\sigma$ maps $\eta+ux+xt$ onto $\sigma(\eta)+uy+yt$, and $x$ is not
a sink, so the values agree after replacing $\eta$ by $\sigma(\eta)$.
In the second line, both sets contain the path $u\to x\to y\to t$, and their remaining
local edges $xt$ and $uy$ are redundant.
As in the NN case, the source and the target then contribute
$(1+q)\,q^{|\eta|+2}$ times the same two values, after reindexing by $\sigma$,
which proves the second relation.
\end{proof}

Multiplying~\eqref{eq:layerGenerating} by $q^{-|S|}$, and using that $\Gamma$ and $R$
have the same number of strict edges, the
NN relation becomes $F(\Gamma)=(1+q)F(R)$ and the DN relation becomes
$F(\Gamma)=F(R)$, for $F(\Gamma)=\LLT_\Gamma(\xvec;q+1)$ and for the right-hand side of
Theorem~\ref{thm:AS}.
These correspond to the shifted forms of the first two bounce relations
of~\cite[Thm.~2.1]{AlexanderssonSulzgruber2022} for a single bounce point.

\begin{remark}\label{rem:fourRelations}
For comparison, we collect the relations
of~\cite[Thm.~2.1 and Thm.~2.5]{AlexanderssonSulzgruber2022} after the
substitution $q\mapsto q+1$.
There $\Gamma$ is encoded by a Schröder path $P$ in the steps
$\nstep,\dstep,\estep$, and we write $F(P)=\LLT_P(\xvec;q+1)$.
Then $F(\nstep\dstep^k\estep)=\elementaryE_{k+1}$, $F$ is multiplicative under
concatenation, and
\begin{align*}
 F(U\nstep\estep V)&=F(U\estep\nstep V)+qF(U\dstep V)
 &&\text{(corner)},\\
 F(U\nstep.\nstep V\dstep.\estep W)&=(1+q)F(U\nstep.\nstep V\estep.\dstep W)
 &&\text{(NN)},\\
 F(U\dstep.\nstep V\dstep.\estep W)&=F(U\nstep.\dstep V\estep.\dstep W)
 &&\text{(DN)},\\
 F(U\nstep.\dstep V\dstep.\estep W)&=qF(U\nstep.\dstep V\estep.\dstep W)
 +(1+q)F(U\dstep.\nstep V\estep.\dstep W)
 &&\text{(ND)},
\end{align*}
where in the last three relations $V\in\{\nstep,\dstep\}^*$ and the bounce path
at the marked points has a single bounce point.
The corner relation deletes a normal corner, or marks it strict; it follows by
summing Lemma~\ref{lem:decoration} over all flip classes, since every corner is
present throughout a flip class or absent from it.
The NN and DN relations are the summed forms of Proposition~\ref{prop:layer}.
We do not use the ND relation: in Lemma~\ref{lem:firstContact}, a strict corner
ending at $y$ is unmarked instead, by Lemma~\ref{lem:decoration}(i).
\end{remark}

\begin{example}\label{ex:NNcontact}
In the graph $\Gamma$ of Example~\ref{ex:running}, with path
$\mathtt{nnnddeneee}$, area sequence $\avec=(0,1,2,3,3,2)$ and strict edges $14$ and
$25$, the triple $(2,3,5)$ is an NN contact: $25$ is the only strict edge leaving
$2$, no strict edge ends at $2$ or $3$, $b(2)=b(3)=1$, and $b(6)=4>3$.
In the notation of Remark~\ref{rem:fourRelations}, the path is
$U\nstep.\nstep V\dstep.\estep W$ with $U=\mathtt{n}$, $V=\mathtt{d}$ and
$W=\mathtt{neee}$.
Its target $R$ has the path $U\nstep.\nstep V\estep.\dstep W=\mathtt{nnndedneee}$,
with area sequence $(0,1,2,3,2,2)$ and strict edges $14$ and $35$:
\[
\LLTpic{nnnddeneee}{}
\;\;\longmapsto\;\;
\LLTpic{nnndedneee}{}
\]
The graph $\Gamma$ has nine normal edges and $R$ has eight, and
Proposition~\ref{prop:layer}, summed as above, gives
$\LLT_\Gamma(\xvec;q+1)=(1+q)\LLT_R(\xvec;q+1)$.
\end{example}

\subsection{Three smaller classes}

We now express certain flip classes in terms of flip classes of graphs with
fewer natural edges.
In all three cases the coefficients are the same for $\Kv$ and $\Ev$, and for
the values with the color fixed at $c=n$.

\begin{lemma}\label{lem:absentCorner}
If some corner is absent from $\theta\cup S$ for $\theta\in C$, then $C$ is a
flip class of a graph with fewer natural edges, with the same values.
\end{lemma}

\begin{proof}
This is Lemma~\ref{lem:decoration}(ii).
\end{proof}

\begin{lemma}\label{lem:disconnected}
Suppose that the natural graph of $\Gamma$ is connected, but the underlying
undirected graph of $\theta\cup S$ is disconnected for $\theta\in C$.
Then there is a graph $\Gamma'$ with fewer natural edges on the same vertex
set, and a flip class $C'$ of $\Gamma'$, with the same values, also with the
color fixed at $c=n$.
\end{lemma}

\begin{proof}
The connected components $B_1,\dotsc,B_r$ of $\theta\cup S$ are the same for
all $\theta\in C$.
List them with the component of $n$ last and the others by their least
elements, and let $\rho$ be the permutation of $[n]$ placing them as consecutive
intervals in this order, keeping the order inside each component.
Let $\Gamma'$ be the disjoint union of the graphs induced by $\Gamma$ on the
components, placed on these intervals.
An induced subgraph of a unit-interval graph on a subset, with the inherited
order, is again a unit-interval graph, and its strict edges are still corners.
Since the natural graph of $\Gamma$ is connected, some natural edge joins two
components, and it is lost; hence $|E(\Gamma')|<|E(\Gamma)|$.

Every edge of $\theta\cup S$, every flip, and every directed path between two
vertices of a component lies inside that component.
Hence $\rho$ maps $C$ bijectively onto a flip class $C'$ of $\Gamma'$, and the
cover tests agree.
Relabelling colorings by $\rho$ preserves $\Kp$.
The highest reachable vertex of $z$ lies in the component of $z$, and $\rho$ is
increasing on each component; so $\rho$ maps fibres to fibres, and fixes $n$
together with the position of $n$ in its fibre.
\end{proof}

\begin{lemma}\label{lem:noCover}
Suppose that every corner is strict, and that for every $\theta\in C$ some corner
is not a cover of $\Pplus{\theta}$.
For a set $T$ of corners, let $\Gamma\setminus T$ be the graph with the corners
in $T$ deleted.
Then the orientations in $C$, viewed as orientations of $\Gamma\setminus T$,
form a union of flip classes of $\Gamma\setminus T$, and
\[
 \Psi(C)=\sum_{\emptyset\neq T}(-1)^{|T|+1}\,\Psi(C;\Gamma\setminus T)
\]
for each of the values $\Psi$ above, where $\Psi(C;\Gamma\setminus T)$ is the value
of $C$ as a set of orientations of $\Gamma\setminus T$, and the sum ranges over
nonempty sets of corners.
\end{lemma}

\begin{proof}
Let $e$ be a natural edge and $f\neq e$ a corner.
A directed path between the endpoints of $e$ stays inside the interval of $e$,
and $f$ is not contained in that interval, since $e$ lies inside some corner and
$f$ is maximal.
So deleting corners other than $e$ does not change whether $e$ is a cover.
The orientations of $\Gamma\setminus T$ are the same sets of ascending edges,
and the flips of $\Gamma\setminus T$ are exactly the flips of $\Gamma$ whose
long edge is not in $T$.
Hence the image of $C$ is closed under flips of $\Gamma\setminus T$, that is,
it is a union of flip classes.

Fix $\theta\in C$ and a corner $f$ which is not a cover of $\Pplus{\theta}$.
The alternative path between the endpoints of $f$ lies strictly inside $f$ and
uses no corner.
It survives the deletion of any set of corners, so by
Lemma~\ref{lem:ineffective}(ii) the sets $T$ and $T\cup\{f\}$ give the same
reachability whenever $f\notin T$.
Pairing these sets gives $\sum_T(-1)^{|T|}\Psi((\theta\cup S)\setminus T)=0$, the
sum ranging over all sets of corners.
The corner $f$ may depend on $\theta$, but the identity holds for each $\theta$.
Summing over $C$ and separating the term $T=\emptyset$ proves the claim.
\end{proof}

\subsection{One flip class}

The following theorem is proved in Section~\ref{sec:connectivity}.
A vertex $v$ is an \defin{articulation} of a connected graph if deleting it
disconnects the graph.

\begin{theorem}\label{thm:connectivity}
Suppose that the natural graph of $\Gamma$ consists of a path $1,2,\dotsc,b$
followed by a connected graph on $[b,n]$ without articulations, where the path
may consist of the single vertex $1=b$.
Let $\theta,\theta'$ be orientations with the same number of ascending edges
such that the underlying graphs of $\theta\cup S$ and $\theta'\cup S$ are
connected and every corner is a cover of both $\Pplus{\theta}$ and
$\Pplus{\theta'}$.
Then $\theta$ and $\theta'$ lie in the same flip class.
\end{theorem}

We shall express flip classes in the form
\begin{equation}\label{eq:reduction}
 \Psi(C)=\sum_iz_i\Psi(C_i)\qquad\text{for }\Psi\in\{\Kv,\Ev,\Kv[\beta],\Ev[\beta]\},
\end{equation}
with the color fixed at $c=n$, where the integers $z_i$ do not depend on $\Psi$,
and every $C_i$ is a flip class of a graph $\Gamma_i$ on $[n]$ such that
$(|E(\Gamma_i)|,\mom(\Gamma_i))$ is lexicographically smaller than
$(|E(\Gamma)|,\mom(\Gamma))$.

\begin{proposition}\label{prop:isolation}
Suppose that $\Gamma$ satisfies the hypothesis of Theorem~\ref{thm:connectivity}
and has an NN or a DN contact.
Let $C$ be a flip class containing an orientation $\theta$ such that
$\theta\cup S$ is connected and every corner is a cover of $\Pplus{\theta}$.
Then $C$ has an expansion~\eqref{eq:reduction}.
\end{proposition}

\begin{proof}
Let $s$ be the number of edges of $\theta\cup S$.
Every other flip class $C'\subseteq\Lay_s(\Gamma)$ omits a corner, or has a
disconnected graph, or contains no orientation in which all corners are covers;
otherwise it would contain an orientation as in Theorem~\ref{thm:connectivity},
and hence equal $C$.
In the first two cases Lemmas~\ref{lem:absentCorner} and~\ref{lem:disconnected}
give an expansion~\eqref{eq:reduction} of $C'$.
In the third case, all corners are present in $C'$, and we mark them strict by
Lemma~\ref{lem:decoration}(i) and apply Lemma~\ref{lem:noCover}.
Proposition~\ref{prop:layer} expresses $\Psi(\Lay_s(\Gamma))$ in terms of one or two
of the sets $\Lay_r(R)$ of the target $R$, which are unions of flip classes of
$R$.
Hence
\[
 \Psi(C)=\Psi(\Lay_s(\Gamma))-\sum_{C'\neq C}\Psi(C')
\]
is an expansion~\eqref{eq:reduction}.
\end{proof}

\subsection{Gluing at one vertex}

Suppose that $c\in[n]$ is such that no natural edge joins a vertex smaller
than $c$ to a vertex larger than $c$.
Let $\Gamma_{\leq c}$ and $\Gamma_{\geq c}$ be the graphs induced on $[1,c]$ and
on $[c,n]$; we say that $\Gamma$ is obtained by \defin{gluing} them at $c$.
An orientation of $\Gamma$ is a pair of orientations of the two parts.
A directed path between two vertices of one part stays in that part, and a
triangle $u<v<w$ with $uw$ natural lies in one part.
Hence the flip classes of $\Gamma$ are the products $C_1\times C_2$ of flip
classes of the two parts.

\begin{figure}[!ht]
\centering
\begin{tikzpicture}[x=2.6em,y=2.6em]
\draw[llt arbitrary,line width=0.6pt] (-0.5,0)--(7.5,0);
\foreach \i in {0,3,5,7} \node[vtx,fill=white] (v\i) at (\i,0) {};
\node[vtx,fibre2] (v1) at (1,0) {};
\node[vtx,fibre2] (v2) at (2,0) {};
\node[vtx,fibre23] (v4) at (4,0) {$c$};
\node[vtx,fibre3] (v6) at (6,0) {};
\draw[sel] (v1) to[bend left=40] (v2);
\draw[sel] (v2) to[bend left=40] (v4);
\draw[sel] (v4) to[bend left=40] (v6);
\draw[decorate,decoration={brace,amplitude=4pt,mirror}] (-0.3,-0.45)--(4.3,-0.45)
  node[midway,below=4pt] {$[1,c]$};
\draw[decorate,decoration={brace,amplitude=4pt}] (3.7,0.5)--(7.3,0.5)
  node[midway,above=4pt] {$[c,n]$};
\end{tikzpicture}
\caption{Gluing at $c$.
A left fibre whose highest reachable vertex is $c$ (hatched) merges with the right
fibre containing $c$ (dotted); $c$ lies in both.}
\label{fig:gluing}
\end{figure}

\begin{lemma}\label{lem:gluing}
Let $C_2$ be a flip class of $\Gamma_{\geq c}$, and suppose that, with the
color fixed at $c$,
\[
 \Psi(C_1)[\beta]=\sum_iz_i\Psi(C_{1,i})[\beta]
 \qquad\text{for all $\beta$ and $\Psi\in\{\Kv,\Ev\}$},
\]
where $C_{1,i}$ is a flip
class of a graph $\Gamma_{1,i}$ on $[1,c]$.
Then $\Psi(C_1\times C_2)=\sum_iz_i\Psi(C_{1,i}\times C_2)$ for $\Psi\in\{\Kv,\Ev\}$, where
$C_{1,i}\times C_2$ is a flip class of the graph obtained by gluing
$\Gamma_{1,i}$ and $\Gamma_{\geq c}$ at $c$.
\end{lemma}

\begin{proof}
The glued graphs are decorated unit-interval graphs: a corner of the left part
ends at most at $c$ and a corner of the right part starts at least at $c$, so
neither contains the other.
For orientations $\theta_1$ and $\theta_2$ of the two parts, a strict
$\Pplus{\theta_1\cup\theta_2}$-partition is a pair of strict partitions of the
parts which agree at $c$.
Hence
\[
 \Kp_{\theta_1\cup\theta_2}=\sum_\beta\Kp_{\theta_1}[\beta]\,
 \widehat{\Kp}_{\theta_2}[\beta],
\]
where $\widehat{\Kp}_{\theta_2}[\beta]$ is the sum over strict partitions of
$\Pplus{\theta_2}$ with value $\beta$ at $c$, with the factor $x_\beta$ omitted.

A left vertex reaches the right part only through $c$.
Hence the highest reachable vertex of a left vertex $u$ is
$\hrv{\theta_1}{u}$ if this is smaller than $c$, and $\hrv{\theta_2}{c}$
otherwise, while the right vertices keep their highest reachable vertices.
The fibre of $c$ in the left part, of size $p$, merges with the fibre of $c$ in
the right part, of size $r$, and all other fibres are unchanged, see
Figure~\ref{fig:gluing}.
In the two chains, $c$ is the largest element of the first and the smallest
element of the second.
Since
\[
 \sum_\beta x_\beta\,\elementaryE_{p-1}(x_1,\dotsc,x_{\beta-1})\,
 \elementaryE_{r-1}(x_{\beta+1},x_{\beta+2},\dotsc)=\elementaryE_{p+r-1}(\xvec),
\]
we obtain $\elementaryE_{\lambda(\theta_1\cup\theta_2)}
=\sum_\beta\elementaryE_{\theta_1}[\beta]\,\widehat{\elementaryE}_{\theta_2}[\beta]$,
with $\widehat{\elementaryE}_{\theta_2}[\beta]$ defined in the same way.
Multiplying the given identity by $\widehat\Psi(C_2)[\beta]$ and summing over
$\beta$ proves the lemma.
\end{proof}

If $\Gamma$ is glued at $c$ from parts with $e_1$ and $e_2$ natural edges, then
$\Gamma$ has $e_1+e_2$ natural edges, and
$\mom(\Gamma)=\mom(\Gamma_{\leq c})+\mom(\Gamma_{\geq c})+(c-1)e_2$, since the
area sequence of the right part starts with $a_c=0$.
Hence replacing the left part by a graph which is smaller in
$(|E|,\mom)$ makes the glued graph smaller as well.

\subsection{The prefix}

Suppose that the natural graph of $\Gamma$ is connected, that is, $a_i>0$ for
$i>1$.
If no $a_i$ exceeds one, the natural graph is the path $1,2,\dotsc,n$.

\begin{lemma}\label{lem:prefix}
Suppose that the natural graph of $\Gamma$ is connected and not a path.
Let $j$ be the least index with $a_j\geq2$, put $b=j-2$, and let $c\geq j$ be the
least index such that $c=n$ or $a_{c+1}=1$.
Then $\Gamma$ is glued at $c$ from $\Gamma_{\leq c}$ and $\Gamma_{\geq c}$, and
the natural graph of $\Gamma_{\leq c}$ is the path $1,\dotsc,b$ followed by a
graph on $[b,c]$ without articulations.
\end{lemma}

\begin{proof}
A vertex $v$ with $1<v<n$ is an articulation of a connected unit-interval graph
if and only if $a_{v+1}=1$.
Indeed, if $a_{v+1}=1$ then $b(v+1)=v$, so no natural edge passes over $v$.
If $a_{v+1}\geq2$, the natural edge $(v-1)(v+1)$ bypasses $v$ in the path of
edges $i(i+1)$, which are natural since all $a_i$ are positive.
By the choice of $j$ we have $a_i=1$ for $1<i<j$ and $a_j=2$, so $b=b(j)$ and
the vertices $2,\dotsc,b$ are articulations.
For $i>b$ we have $a_i\leq i-b$, so no natural edge ends above $b$ and starts
below $b$.
By the choice of $c$, no vertex strictly between $b$ and $c$ is an
articulation of the graph induced on $[b,c]$, and $b(c+1)=c$ if $c<n$.
\end{proof}

\begin{figure}[!ht]
\centering
\LLTpic[cell=0.9em]{nnennneeneenneenee}{\LLTshade{1/2}{part1}\LLTshade{2/3,2/4,3/4,2/5,3/5,4/5,4/6,5/6}{part2}\LLTshade{6/7,6/8,7/8,8/9}{part3}}
\caption{The area sequence $\avec=(0,1,1,2,3,2,1,2,1)$, with $j=4$, $b=2$ and $c=6$ in
Lemma~\ref{lem:prefix}.
The prefix is the path $1,2$ (dotted) followed by the graph on $[2,6]$ without
articulations (hatched), and the whole graph is glued at $6$ from the prefix and
the graph on $[6,9]$ (cross-hatched).}
\label{fig:prefix}
\end{figure}

We call $\Gamma_{\leq c}$ the \defin{prefix} of $\Gamma$, see Figure~\ref{fig:prefix}.
For the prefix we put $a_{c+1}=0$.
The prefix has the same area sequence as $\Gamma$ up to $c$, and
$a_{c+1}<a_c$ both in $\Gamma$ and in the prefix, since $a_c\geq2$.

\begin{lemma}\label{lem:firstContact}
Let $\Gamma$ be a graph in which every corner is strict, with connected natural
graph which is not a path.
Let $t$ be the least index with $a_{t+1}<a_t$, where $a_{n+1}=0$, and put
$x=b(t)$ and $y=x+1$.
Then $a_t\geq2$, and after unmarking the strict edge ending at $y$, if there is
one, $(x,y,t)$ is an NN contact or $(b(x),x,y,t)$ is a DN contact.
\end{lemma}

For the graph of Example~\ref{ex:NNcontact}, $t=5$ and $(2,3,5)$ is an NN contact;
the two cases are drawn in Figure~\ref{fig:contacts}.

\begin{proof}
The areas $a_1,\dotsc,a_t$ are weakly increasing.
If $a_t=1$, then all these areas are at most one, and since $a_{t+1}=0$, the
natural graph would be a path; hence $a_t\geq2$ and $x<y<t$.
Since $a_{t+1}<a_t$, the edge $xt$ is a corner, hence strict, and
$b(t+1)\geq x+2>y$.

For $i<t$ we have $b(i+1)-b(i)\in\{0,1\}$, and a corner ends at $i$ exactly when
$b(i+1)=b(i)+1$.
Hence the corners with heads at most $t$ have the tails $1,2,\dotsc,x$, each
exactly once, and the corners with heads above $t$ have tails at least
$b(t+1)\geq x+2$.
In particular, no strict edge leaves $y$, and $xt$ is the only strict edge
leaving $x$.

Since $y<t$, we have $a_x\leq a_y\leq a_x+1$.
If $a_y=a_x+1$, then $b(x)=b(y)$, no corner ends at $x$, and $(x,y,t)$ is an NN
contact once the corner ending at $y$, if any, is unmarked.
If $a_y=a_x$, then $b(y)=b(x)+1$.
The pair $xy$ is natural, so $a_x=a_y\geq1$, and $u=b(x)$ is the tail of the
strict corner $ux$.
Then $(u,x,y,t)$ is a DN contact once the corner ending at $y$, if any, is
unmarked.
\end{proof}

In the terminology of~\cite{AlexanderssonSulzgruber2022}, the four cases are
the contacts NN, ND, DN and DD at the first bounce point.
By Lemma~\ref{lem:decoration}(i), unmarking the corner at $y$ in the ND and DD
cases does not change a flip class in which this corner is present.

\subsection{Proof of Theorem~\ref{thm:main} for ordinary graphs}

We proceed by induction on $n$, then on $|E(\Gamma)|$, then on $\mom(\Gamma)$.
For $n\leq1$ there is only one orientation, and both values are $1$ or
$\elementaryE_1$.

Let $C$ be a flip class of $\Gamma$.
If some corner is absent from $C$, we apply Lemma~\ref{lem:absentCorner} and
induction.
Otherwise we mark all corners strict, by Lemma~\ref{lem:decoration}(i).
If the natural graph is disconnected, it is a disjoint union of graphs on
consecutive intervals; the flip class is a product of flip classes of these
graphs, both values are products, and we use induction on $n$.
If the natural graph is a path, all its edges are strict corners, there is
only one orientation, and both values equal $\elementaryE_n$.
If $\theta\cup S$ is disconnected for $\theta\in C$, or no orientation in $C$
has all corners as covers, we apply Lemma~\ref{lem:disconnected}
or~\ref{lem:noCover} and induction.

In the remaining case, let $\Gamma$ be glued at $c$ from its prefix and
$\Gamma_{\geq c}$ as in Lemma~\ref{lem:prefix}, and write $C=C_1\times C_2$.
Let $\theta\in C$ be connected with all corners as covers.
Its restriction to the prefix is connected, and every corner of the prefix is a
corner of $\Gamma$ ending at most at $c$, whose cover test only involves
vertices up to $c$.
By Lemma~\ref{lem:firstContact} and Lemma~\ref{lem:decoration}(i), the prefix
has an NN or DN contact, possibly after unmarking one corner which is present in
$C_1$.
The prefix satisfies the hypothesis of Theorem~\ref{thm:connectivity}, and
Proposition~\ref{prop:isolation} gives an expansion~\eqref{eq:reduction}
$\Psi(C_1)[\beta]=\sum_iz_i\Psi(C_{1,i})[\beta]$ on $[1,c]$, with the color fixed at
the last vertex $c$ of the prefix.
By Lemma~\ref{lem:gluing},
\[
 \Psi(C)=\sum_iz_i\Psi(C_{1,i}\times C_2),\qquad \Psi\in\{\Kv,\Ev\}.
\]
Each $C_{1,i}\times C_2$ is a flip class of a graph on $[n]$ which is smaller in
$(|E|,\mom)$, by the remark after Lemma~\ref{lem:gluing}.
By induction $\Kv(C_{1,i}\times C_2)=\Ev(C_{1,i}\times C_2)$ for all $i$, and
hence $\Kv(C)=\Ev(C)$.
This completes the proof of Theorem~\ref{thm:main} for ordinary graphs, given
Theorem~\ref{thm:connectivity}.
